\appendix
\section{Zorn's lemma}

\begin{de}
A \emph{partial order} on a set $S$ is a relation $\leq$ on $S$ satisfying,
for all $x,y,z\in S$:

(a) $x\leq x$ (reflexivity);

(b) if $x\leq y$ and $y\leq x$, then $x=y$ (antisymmetry);

(c) if $x\leq y$ and $y\leq z$, then $x\leq z$ (transitivity).

The pair $(S,\leq)$ is called a \emph{partially ordered set}, or a
\emph{poset}.

Let $T\subseteq S$. An element $u\in S$ is an \emph{upper bound} for $T$ if
$t\leq u$ for every $t\in T$; the upper bound need not belong to $T$.
An element $m\in T$ is a \emph{maximal element} of $T$ if, whenever
$m\leq t$ for some $t\in T$, one has $t=m$.

The poset $(S,\leq)$ is \emph{linearly ordered} (or \emph{totally ordered})
if, for every $x,y\in S$, either $x\leq y$ or $y\leq x$. A \emph{chain} in
$S$ is a subset of $S$ that is linearly ordered by the restriction of
$\leq$.
\end{de}

\begin{thm}[Zorn's lemma]
Let $(S,\leq)$ be a nonempty partially ordered set. If every chain in $S$
has an upper bound in $S$, then $S$ has a maximal element.
\end{thm}

\begin{thm}
Every vector space over $\mathbb{R}$ has a basis.
\end{thm}

\begin{proof}
Let $V$ be a vector space over $\mathbb{R}$, and let
\[
\mathcal{P}=\{A\subseteq V:A\text{ is linearly independent}\},
\]
ordered by inclusion. The set $\mathcal{P}$ is nonempty because
$\varnothing\in\mathcal{P}$.

Let $\mathcal{C}$ be a chain in $\mathcal{P}$, and set
\[
B=\bigcup_{A\in\mathcal{C}}A.
\]
We claim that $B$ is linearly independent. Suppose
$x_1,\ldots,x_n\in B$ are distinct and
\[
a_1x_1+\cdots+a_nx_n=0.
\]
For each $i$, choose $A_i\in\mathcal{C}$ such that $x_i\in A_i$. Because
$A_1,\ldots,A_n$ belong to a chain, one of these finitely many sets contains
all the others. Hence all the vectors $x_1,\ldots,x_n$ belong to one
linearly independent set $A_k$, and therefore
\[
a_1=\cdots=a_n=0.
\]
Thus $B$ is linearly independent, so $B\in\mathcal{P}$. Moreover, every
$A\in\mathcal{C}$ satisfies $A\subseteq B$, so $B$ is an upper bound for
$\mathcal{C}$ in $\mathcal{P}$. This also covers the empty chain, for which
$B=\varnothing$.

By Zorn's lemma, $\mathcal{P}$ has a maximal element $M$. It remains to show
that $M$ spans $V$. Suppose, to the contrary, that there is a vector
$v\in V\setminus\operatorname{span}M$. We claim that $M\cup\{v\}$ is
linearly independent. Suppose $m_1,\ldots,m_n\in M$ are distinct and
\[
av+b_1m_1+\cdots+b_nm_n=0.
\]
If $a\neq0$, then
\[
v=-a^{-1}(b_1m_1+\cdots+b_nm_n)\in\operatorname{span}M,
\]
a contradiction. Hence $a=0$, and the linear independence of $M$ gives
$b_1=\cdots=b_n=0$. Thus $M\cup\{v\}$ is linearly independent and properly
contains $M$, contradicting the maximality of $M$.

Therefore $\operatorname{span}M=V$, so $M$ is a basis of $V$.
\end{proof}